% Appendices

\chapter{Experimental Configuration and Parameters}

\section{Simulation Parameters}

Table \ref{tab:sim_params} summarizes the key simulation parameters used in all experiments:

\begin{table}[h]
\centering
\small
\caption{Simulation configuration parameters}
\label{tab:sim_params}
\begin{tabular}{|l|l|}
\hline
\textbf{Parameter} & \textbf{Value} \\
\hline
Grid dimensions & $30 \times 30$ cells (900 total) \\
Grid topology & Non-toroidal (no wrapping) \\
Time horizon & 500 discrete time steps \\
Agent movement & Moore neighborhood (8 directions) \\
Distance metric & Chebyshev distance ($L_\infty$) \\
\hline
\end{tabular}
\end{table}

\section{Experimental Design Matrix}

Table \ref{tab:design_matrix} presents the full factorial design:

\begin{table}[h]
\centering
\small
\caption{Full factorial experimental design (20 configurations × 20 replications = 400 runs)}
\label{tab:design_matrix}
\begin{tabular}{|c|c|c|c|}
\hline
\textbf{Communication Range} & \textbf{Team Size} & \textbf{Resource Density} & \textbf{Runs} \\
\hline
$r \in \{0, 2, 4, 6, 8\}$ & $N \in \{10, 20\}$ & $\rho \in \{0.15, 0.25\}$ & 20 per config \\
\hline
\end{tabular}
\end{table}

\section{BDI Agent Parameters}

\begin{itemize}
    \item \textbf{Local sensing radius:} 2 cells (Moore neighborhood)
    \item \textbf{Movement strategy:} Greedy (minimize Chebyshev distance to target)
    \item \textbf{Exploration strategy:} Random movement to adjacent cells
    \item \textbf{Belief update:} Set union of received messages
    \item \textbf{Target selection:} Nearest known fruit location
    \item \textbf{Idle threshold:} 2 consecutive steps without finding fruit
\end{itemize}

\section{Communication Protocol Details}

\subsection{Message Format}

Each message contains:
\begin{itemize}
    \item Sender agent ID
    \item Sender position $(x, y)$
    \item Belief set: list of known fruit locations
    \item Timestamp (simulation step)
\end{itemize}

\subsection{Chebyshev Distance Neighborhoods}

For each communication range $r$, the neighborhood size is $(2r+1) \times (2r+1)$:

\begin{table}[h]
\centering
\small
\caption{Communication neighborhood sizes by range}
\label{tab:neighborhoods}
\begin{tabular}{|c|c|c|}
\hline
\textbf{Range} & \textbf{Neighborhood Size} & \textbf{Max Agents} \\
\hline
0 & N/A & 0 \\
2 & $5 \times 5$ & 24 \\
4 & $9 \times 9$ & 80 \\
6 & $13 \times 13$ & 168 \\
8 & $17 \times 17$ & 288 \\
\hline
\end{tabular}
\end{table}

\section{Data Collection Metrics}

\subsection{Primary Metrics}

\begin{itemize}
    \item \textbf{Total Yield ($Y$):} Cumulative fruit harvested by all agents
    \item \textbf{Total Messages ($M$):} Cumulative messages sent during episode
    \item \textbf{Remaining Fruit:} Unharvested fruit at episode end
\end{itemize}

\subsection{Derived Metrics}

\begin{itemize}
    \item \textbf{Communication Efficiency:} $E = Y / M$ (yield per message)
    \item \textbf{Harvest Rate:} Fraction of available fruit harvested
    \item \textbf{Messages per Step:} $M / T$ (average communication overhead)
\end{itemize}

\section{Reproducibility Information}

\subsection{Software Environment}

\begin{itemize}
    \item \textbf{Framework:} Mesa 3.0 (Python agent-based modeling)
    \item \textbf{Python version:} 3.9+
    \item \textbf{Key dependencies:} numpy, pandas, matplotlib
    \item \textbf{Random seed range:} 0-19 (fixed for reproducibility)
\end{itemize}

\subsection{Repository Information}

Complete source code, experimental data, and analysis scripts are available at:

\begin{center}
\url{https://github.com/marufdevops/multiagent-cooperation-research}
\end{center}

This repository includes:
\begin{itemize}
    \item Complete simulation implementation
    \item Batch experiment runner
    \item Raw experimental results (CSV format)
    \item Analysis and visualization scripts
    \item This dissertation source (LaTeX)
\end{itemize}

